\documentclass[11pt]{article}
\usepackage[margin=1in]{geometry}
\usepackage{amsmath,amssymb,booktabs,graphicx,hyperref,natbib}
\title{From CLV Ranking to Uplift-Aware Targeting:\\
Cross-Context Evidence from Online Retail, Grocery, and Supermarket Campaign Data}
\author{Group 11}
\date{}

\begin{document}
\maketitle

\begin{abstract}
Customer lifetime value (CLV) models identify \emph{who is valuable}; uplift models identify
\emph{who is responsive} to marketing. We connect the two across three retail contexts. On
UK online retail and US grocery data we show an explainable Hurdle--Semantic CLV framework
transfers across contexts but for \emph{different reasons}: in zero-inflated e-commerce the
return-incidence stage dominates, while in habitual grocery shopping temporal-sequence signal
dominates. On the X5 RetailHero campaign data (randomised treatment/control), we show that
ranking customers by predicted response is \emph{worse than random} for incremental targeting,
whereas uplift learners recover persuadable customers. Finally, a \emph{value-adjusted uplift}
score ($\hat\tau_i\cdot\hat V_i$) yields the highest simulated profit under tight budgets,
though we report transparently that the dollar differences are directional rather than
statistically significant. CLV tells us who is valuable; uplift tells us who is responsive;
value-adjusted uplift tells us who is worth targeting.
\end{abstract}

\section{Introduction}
CLV studies typically rank high-value customers, while uplift studies rank campaign-responsive
customers. Yet a high-CLV customer may purchase regardless of intervention, and a mid-value but
high-uplift customer may be the true ``persuadable''. Limited work systematically connects
\emph{explainable CLV ranking} with \emph{uplift-aware targeting} across multiple retail contexts.
We study four questions:
\begin{itemize}
  \item[\textbf{RQ1}] How robust is the explainable Hurdle--Semantic CLV framework across UK
  online retail and US grocery contexts?
  \item[\textbf{RQ2}] Do semantic product profiles and behavioural sequence features add
  incremental value beyond RFM for CLV ranking, calibration, and interpretation?
  \item[\textbf{RQ3}] On campaign data with treatment/control labels, does uplift-aware targeting
  outperform CLV-only, RFM-only, and response-model targeting?
  \item[\textbf{RQ4}] Can combining predicted value with uplift improve simulated campaign
  efficiency, and what transferable principles follow for budget-constrained markets?
\end{itemize}

\section{Related Work}
Zero-inflated / heavy-tailed CLV motivates Hurdle and ZILN decompositions \citep{cragg1971hurdle,
wang2019deep}; sequence features improve CLV \citep{bauer2021improved}; customer embeddings give
semantic representations \citep{chamberlain2017customer}. Conformalized quantile regression gives
valid prediction intervals \citep{romano2019cqr}. Uplift modelling estimates incremental treatment
effect \citep{gutierrez2017uplift}; meta-learners and causal forests estimate heterogeneous
effects \citep{kunzel2019metalearners,wager2018causalforest}; causal ML supports campaign
targeting \citep{gubela2024multiple,fader2005rfm}.

\section{Data and Problem Setup}
\textbf{Online Retail II} (UK e-commerce, Phase~1 baseline, read-only) and \textbf{Dunnhumby
Complete Journey} (US grocery, external CLV validation) provide CLV ranking; \textbf{X5
RetailHero} (randomised SMS campaign) provides treatment/control for uplift. We deliberately do
\emph{not} use Dunnhumby coupon redemption as a causal treatment (it is a post-campaign outcome,
not clean assignment). CLV target $Y_i=\sum_{t\in\text{horizon}}\text{Revenue}_{i,t}$.

\section{Methodology}
\textbf{Module A --- CLV ranking (Hurdle).} $\widehat{CLV}_i=p_i\cdot m_i$ with
$p_i=P(Y_i>0\mid X_i)$ and $m_i=E[Y_i\mid Y_i>0,X_i]$, using a lognormal bias correction.
\textbf{Module B --- Semantic profiling.} Product co-purchase PPMI $\to$ truncated SVD
embeddings; customer taste vector = revenue-weighted mean of bought-product embeddings.
\textbf{Module C --- Uplift.} $\tau_i=E[Y_i(1)-Y_i(0)\mid X_i]$ via S/T/X-Learners and class
transformation. \textbf{Module D --- Value-adjusted uplift.}
$\text{Score}^{\text{VU}}_i=\hat\tau_i\cdot\hat V_i$, with $\hat V_i$ a value proxy from
historical spend (\emph{not} actual future CLV).

\section{Experiments and Results}
We use a stratified 80/20 split (CLV) and 70/30 split (uplift), seed 42. Metrics: Normalized
Gini (NG), Revenue Capture@10\% (RC@10), Precision@10\%, decile calibration; Qini AUC, AUUC,
uplift@K; CQR coverage; bootstrap CIs.

\begin{table}[t]\centering
\caption{Walk-forward CLV benchmark. NG and RC@10 are mean $\pm$ std across temporal windows (Online Retail II on the Phase~1 windows W1--W3, Dunnhumby on D1--D4); Precision@10 is the across-window mean. Results are temporally valid and lie in the same range as the Phase~1 walk-forward benchmark (NG~$\approx$~0.83); small differences follow from the feature set and the window subset.}\label{tab:clvwf}
\begin{tabular}{lrrrr}
\hline
context & model & NG & RC@10 & Prec@10 \\
\hline
Online Retail II (W1-W3) & Hurdle & 0.825 $\pm$ 0.050 & 59.037 $\pm$ 9.466 & 0.577 \\
Online Retail II (W1-W3) & Monetary & 0.822 $\pm$ 0.058 & 60.017 $\pm$ 8.404 & 0.586 \\
Online Retail II (W1-W3) & RandomForest & 0.805 $\pm$ 0.043 & 56.883 $\pm$ 7.203 & 0.558 \\
Online Retail II (W1-W3) & Ridge & 0.805 $\pm$ 0.038 & 56.193 $\pm$ 8.234 & 0.547 \\
Online Retail II (W1-W3) & XGBoost & 0.802 $\pm$ 0.044 & 56.807 $\pm$ 8.818 & 0.540 \\
Online Retail II (W1-W3) & LightGBM & 0.799 $\pm$ 0.052 & 57.170 $\pm$ 9.850 & 0.547 \\
Dunnhumby (D1-D4) & RandomForest & 0.862 $\pm$ 0.012 & 34.297 $\pm$ 0.623 & 0.723 \\
Dunnhumby (D1-D4) & Ridge & 0.858 $\pm$ 0.009 & 34.280 $\pm$ 0.452 & 0.703 \\
Dunnhumby (D1-D4) & Hurdle & 0.855 $\pm$ 0.014 & 34.280 $\pm$ 0.457 & 0.714 \\
Dunnhumby (D1-D4) & XGBoost & 0.854 $\pm$ 0.014 & 34.235 $\pm$ 0.923 & 0.738 \\
Dunnhumby (D1-D4) & LightGBM & 0.849 $\pm$ 0.009 & 34.270 $\pm$ 0.778 & 0.723 \\
Dunnhumby (D1-D4) & Monetary & 0.825 $\pm$ 0.010 & 32.350 $\pm$ 0.712 & 0.658 \\
\hline
\end{tabular}
\end{table}

\begin{table}[t]\centering
\caption{X5 uplift model comparison (RQ3). The response model ranks below random.}\label{tab:uplift}
\begin{tabular}{lrrrr}
\hline
model & Qini\_AUC & AUUC & uplift@10 & uplift@20 \\
\hline
S-Learner & 0.014 & 0.020 & 0.112 & 0.077 \\
X-Learner & 0.012 & 0.018 & 0.104 & 0.072 \\
T-Learner & 0.010 & 0.015 & 0.087 & 0.069 \\
ClassTransform & 0.008 & 0.011 & 0.086 & 0.064 \\
Random & -0.002 & -0.004 & 0.026 & 0.029 \\
ResponseModel & -0.007 & -0.011 & 0.004 & 0.008 \\
\hline
\end{tabular}
\end{table}

\begin{table}[t]\centering
\caption{Targeting policy comparison on X5 (RQ4). The value proxy $\hat V_i$ is pre-period historical monetary spend (a gross-value proxy, not future CLV). Incremental conversions and incremental value are held-out treated-minus-control estimates within each selected top-$K$ set (Eq.~\ref{eq:profit}); the headline uses margin $m=1$, so profit $=$ (incr.\ value) $-$ contact cost with $c=100$ rubles per contact (value and profit in millions of rubles). Across a sensitivity sweep ($m\in\{0.1,0.2,0.3\}$, $c\in\{50,100,200\}$) value-adjusted uplift has the highest point-estimate profit in all nine settings.}\label{tab:policy}
\begin{tabular}{lrrrr}
\hline
policy & K & incr. conv. & incr. value (M) & profit (M) \\
\hline
Random & 10\% & 129.4 & 0.29 & -0.31 \\
Random & 20\% & 351.5 & 0.13 & -1.07 \\
RFM-only & 10\% & 52.0 & 0.11 & -0.49 \\
RFM-only & 20\% & 109.7 & 1.96 & 0.76 \\
Value-only(CLV proxy) & 10\% & 77.3 & -0.05 & -0.65 \\
Value-only(CLV proxy) & 20\% & 263.0 & 4.56 & 3.36 \\
Uplift-only(S-Learner) & 10\% & 670.4 & 1.09 & 0.49 \\
Uplift-only(S-Learner) & 20\% & 930.1 & 3.23 & 2.03 \\
Value-adjusted(uplift x value) & 10\% & 248.4 & 1.34 & 0.74 \\
Value-adjusted(uplift x value) & 20\% & 451.7 & 5.23 & 4.03 \\
\hline
\end{tabular}
\end{table}

\begin{table}[t]\centering
\caption{Significance of key claims ($\Delta$NG with walk-forward paired $t$-tests for CLV; bootstrap for the uplift Qini comparison). With only 3--4 windows, the CLV tests are robustness diagnostics rather than definitive inference. Only sequence-in-grocery and uplift-vs-response are significant.}\label{tab:robust}
\begin{tabular}{lrr}
\hline
claim & $\Delta$NG & $p$ \\
\hline
Hurdle $>$ Monetary (e-com, walk-fwd) & $+0.0030$ & $0.554$ \\
Hurdle vs RandomForest (grocery, walk-fwd) & $-0.0064$ & $0.063$ \\
extra features $>$ RFM (e-com) & $+0.0092$ & $0.1668$ \\
extra features $>$ RFM (grocery) & $+0.0444$ & $0.0021$ \\
semantic adds ranking (e-com) & $+0.0028$ & $0.5835$ \\
semantic adds ranking (grocery) & $+0.0011$ & $0.8373$ \\
sequence adds ranking (grocery) & $+0.0206$ & $0.0011$ \\
S-Learner $>$ Response (uplift Qini) & -- & $<0.001$ \\
\hline
\end{tabular}
\end{table}

\begin{table}[t]\centering
\caption{Leave-one-group-out ablation (Hurdle, walk-forward). dNG\_vs\_ALL $>$ 0 means dropping the group hurts; paired\_p is the paired t-test across windows.}\label{tab:ablation}
\begin{tabular}{lrrrrr}
\hline
context & variant & NG\_mean & NG\_std & dNG\_vs\_ALL & paired\_p \\
\hline
Online Retail II & ALL & 0.829 & 0.056 & 0.000 & -- \\
Online Retail II & ALL-behavioural & 0.829 & 0.052 & -0.000 & 0.954 \\
Online Retail II & ALL-semantic & 0.826 & 0.048 & 0.003 & 0.584 \\
Online Retail II & RFM-only & 0.820 & 0.063 & 0.009 & 0.167 \\
Dunnhumby & ALL & 0.856 & 0.015 & 0.000 & -- \\
Dunnhumby & ALL-behavioural & 0.856 & 0.015 & 0.000 & 0.088 \\
Dunnhumby & ALL-semantic & 0.855 & 0.016 & 0.001 & 0.837 \\
Dunnhumby & ALL-sequence & 0.836 & 0.017 & 0.021 & 0.001 \\
Dunnhumby & RFM-only & 0.812 & 0.013 & 0.044 & 0.002 \\
\hline
\end{tabular}
\end{table}

\begin{table}[t]\centering
\caption{Top-10\% targeting overlap between policies, defined as $|A_K\cap B_K|/K$ where $A_K,B_K$ are the top-$K$ customer sets selected by each policy. Value and Uplift select almost disjoint customers (0.007), evidencing that value $\neq$ persuadability.}\label{tab:overlap}
\begin{tabular}{lrr}
\hline
policy\_A & policy\_B & overlap \\
\hline
RFM & Value(CLV proxy) & 0.666 \\
RFM & Uplift(S-Learner) & 0.003 \\
RFM & Value-adjusted & 0.173 \\
Value(CLV proxy) & Uplift(S-Learner) & 0.007 \\
Value(CLV proxy) & Value-adjusted & 0.312 \\
Uplift(S-Learner) & Value-adjusted & 0.256 \\
\hline
\end{tabular}
\end{table}

\begin{table}[t]\centering
\caption{Conformalized quantile regression: empirical coverage and mean interval width per context. Coverage is close to nominal; grocery intervals are tighter, consistent with its lighter tail.}\label{tab:cqr}
\begin{tabular}{lrrr}
\hline
context & nominal\_coverage & empirical\_coverage & mean\_interval\_width \\
\hline
Online Retail II & 0.900 & 0.902 & 1,653 \\
Online Retail II & 0.800 & 0.793 & 1,315 \\
Dunnhumby & 0.900 & 0.890 & 707 \\
Dunnhumby & 0.800 & 0.778 & 515 \\
\hline
\end{tabular}
\end{table}


\textbf{RQ1/RQ2.} Table~\ref{tab:clv_summary}: Hurdle is best on e-commerce (NG$=0.91$) but ties
a linear model on grocery (NG$=0.85$); extra features beat RFM significantly in both contexts
(Table~\ref{tab:robust}), with a larger gain in grocery. SHAP shows monetary drives spend in
e-commerce while recent-sequence spend drives it in grocery; CQR attains valid coverage
(90\%$\to$90.2\%/89.0\%).

\textbf{RQ3.} Table~\ref{tab:uplift}: the response model ranks \emph{below random} (Qini
$-0.007$, CI entirely negative), while uplift learners reach $\approx 11\%$ uplift in the top
decile --- ``likely to buy'' $\neq$ ``persuadable''.

\textbf{RQ4.} Table~\ref{tab:policy}: targeting by value alone can lose money at small budgets;
uplift-only maximises incremental conversions; value-adjusted uplift gives the highest simulated
profit at tight budgets. Bootstrap CIs (Table~\ref{tab:robust}) confirm the ranking claims are
significant, but the \emph{dollar} profit differences are directional, not significant.

\section{Discussion: Transferable Insights}
(1) Do not target on CLV alone --- it selects ``sure buyers''. (2) Do not target on response
probability alone --- it selects low-incremental customers. (3) Value-adjusted uplift fits
budget-constrained settings. (4) Explainability clarifies \emph{why} a customer is selected.
(5) Uncertainty matters when market data are unstable. These provide transferable design
principles for budget-constrained emerging retail markets.

\section{Limitations}
Single train/test split per context (mitigated by bootstrap); X5 value proxy is historical, not
future, spend; profit differences between policies are not statistically significant; uplift
absolute Qini is modest (the campaign signal itself is small, $\approx 3.3$pp).

\section{Conclusion}
Across three retail contexts we connect explainable CLV ranking with uplift-aware targeting. The
framework transfers but its mechanism is context-dependent, and value-adjusted uplift is the most
promising targeting rule under budget constraints.

\bibliographystyle{plainnat}
\bibliography{references}
\end{document}
